\documentclass[11pt,a4paper]{article}
\usepackage[T1]{fontenc}
\usepackage{lmodern}
\usepackage[top=2.2cm,bottom=2.4cm,left=2.3cm,right=2.3cm]{geometry}
\usepackage{microtype}
\usepackage{xcolor}
\usepackage{enumitem}
\usepackage[nobottomtitles*]{titlesec}
\usepackage{etoolbox}
\usepackage{fancyhdr}
\usepackage{lastpage}

% Publications come from refs.bib. biblatex with Biber is the right tool here:
% it can split one .bib file into sections by entry type and sort newest first.
\usepackage[backend=biber,style=numeric,sorting=ydnt,giveninits=true,maxbibnames=99,
            defernumbers=true,url=false,isbn=false]{biblatex}
\addbibresource{refs.bib}

\usepackage[hidelinks]{hyperref}

\definecolor{accent}{HTML}{1F4E79}

% ---- Your name -------------------------------------------------------------------
\newcommand*{\myname}{Alex Morgan}   % shown in the page footer
% In the publication lists, every author whose family name and first initial match
% the next two lines is set in bold. Use the same spelling as in refs.bib ("Morgan, Alex").
\newcommand*{\myfamily}{Morgan}
\newcommand*{\myinitial}{A}

\def\cvfirstletter#1#2\relax{#1}
\newcommand*{\ifmyself}[2]{%
  \ifdefempty{\namepartgiven}{\def\cvgiven{}}{\edef\cvgiven{\expandafter\cvfirstletter\namepartgiven\relax}}%
  \ifthenelse{\equal{\namepartfamily}{\myfamily}\AND\equal{\cvgiven}{\myinitial}}{#1}{#2}}
\DeclareNameFormat{default}{%
  \nameparts{#1}%
  \begingroup
  \ifmyself{\let\mkbibnamefamily\mkbibbold \let\mkbibnamegiven\mkbibbold
            \let\mkbibnameprefix\mkbibbold \let\mkbibnamesuffix\mkbibbold}{}%
  \ifgiveninits
    {\usebibmacro{name:given-family}{\namepartfamily}{\namepartgiveni}{\namepartprefix}{\namepartsuffix}}
    {\usebibmacro{name:given-family}{\namepartfamily}{\namepartgiven}{\namepartprefix}{\namepartsuffix}}%
  \endgroup
  \usebibmacro{name:andothers}}

% ---- Look: headings, entry lists, page footer ---------------------------------
\setlength{\parindent}{0pt}
\setlength{\parskip}{0pt}
\titleformat{\section}{\large\bfseries\color{accent}}{}{0pt}{\MakeUppercase}[\vspace{-0.6ex}\color{accent}\titlerule]
\titlespacing*{\section}{0pt}{1.6em}{0.7em}
\titleformat{\subsection}{\normalsize\bfseries}{}{0pt}{}
\titlespacing*{\subsection}{0pt}{1em}{0.4em}

% One entry per \item: the date goes in the brackets, the rest is free text.
% A long entry can break across pages, and you add or remove items without
% touching any layout.
\newlist{cvlist}{description}{1}
\setlist[cvlist]{leftmargin=3.1cm,labelwidth=2.8cm,labelsep=0.3cm,align=left,
                 font=\normalfont\color{black!70},itemsep=0.55em,topsep=0pt,parsep=0pt}
\setlist[itemize]{nosep,leftmargin=1.2em,label={\textbullet}}

\pagestyle{fancy}
\fancyhf{}
\renewcommand{\headrulewidth}{0pt}
\fancyfoot[L]{\footnotesize\color{black!60} \myname{} \textperiodcentered\ Curriculum vitae \textperiodcentered\ updated \today}
\fancyfoot[R]{\footnotesize\color{black!60} Page \thepage\ of \pageref{LastPage}}

% Publication list: small type, a gap between entries, and a hanging label.
\renewcommand*{\bibfont}{\small}
\setlength{\bibitemsep}{0.35\baselineskip}
\defbibheading{pubs}[\bibname]{\subsection{#1}}

\begin{document}

% ---- Header ---------------------------------------------------------------------
{\Huge\bfseries\color{accent} \myname}\par\smallskip
{\large Assistant Professor of Statistics}\par
Department of Mathematics and Statistics, Example University\par\medskip
\small
Office 4.12, 10 University Road, Springfield \quad
\href{mailto:alex.morgan@example.edu}{alex.morgan@example.edu} \quad
\href{https://example.edu/~amorgan}{example.edu/\textasciitilde amorgan}\\
ORCID \href{https://orcid.org/0000-0000-0000-0000}{0000-0000-0000-0000}
\normalsize

\section{Research interests}
Statistical learning on networks, robust and high-dimensional inference, and Monte Carlo
methods. Two or three lines are enough; the publication list says the rest.

\section{Education}
\begin{cvlist}
  \item[2015--2019] \textbf{PhD in Statistics}, Example Institute of Technology\\
    Thesis: \emph{Sparse Graph Estimation in High Dimensions}. Advisor: Prof.\ Jordan Ellis.
    Awarded with distinction.
  \item[2013--2015] \textbf{MSc in Applied Mathematics}, University of Northfield\\
    Graduated first in a class of 64.
  \item[2010--2013] \textbf{BSc in Mathematics}, University of Northfield
\end{cvlist}

\section{Academic positions}
\begin{cvlist}
  \item[2023--present] \textbf{Assistant Professor of Statistics}, Example University\\
    Tenure-track; leads the Network Inference Group (two PhD students, one postdoc).
  \item[2021--2023] \textbf{Postdoctoral Researcher}, Institute for Data Science, Riverside University\\
    Funded by a Marie Example Fellowship; host: Prof.\ Dana Whitfield.
  \item[2019--2021] \textbf{Postdoctoral Associate}, Example Institute of Technology
  \item[Spring 2018] \textbf{Visiting Researcher}, Department of Statistics, Lakeshore University
\end{cvlist}

\section{Grants and awards}
\subsection{Grants}
\begin{cvlist}
  \item[2025--2028] \textbf{Robust Learning on Networks}, National Science Foundation, \$480{,}000\\
    Principal investigator.
  \item[2024--2025] \textbf{Seed Grant for Early-Career Faculty}, Example University, \$30{,}000
  \item[2021--2023] \textbf{Marie Example Fellowship}, \texteuro{}184{,}000\\
    Sole applicant; success rate in the call 14\%.
\end{cvlist}
\subsection{Awards}
\begin{cvlist}
  \item[2025] Best Paper Award, International Conference on Artificial Intelligence and Statistics
  \item[2020] Dissertation Prize, Example Statistical Society
  \item[2019] Outstanding Teaching Assistant Award, Example Institute of Technology
\end{cvlist}

\section{Teaching}
\subsection{Courses taught}
\begin{cvlist}
  \item[2023--present] \textbf{Statistical Learning} (graduate, 45 students) and
    \textbf{Probability and Statistics} (undergraduate, 120 students), Example University.
    Average student rating 4.6 of 5.
  \item[2022] \textbf{Introduction to Monte Carlo Methods}, guest lecturer, 6 hours, Riverside University
\end{cvlist}
\subsection{Supervision}
\begin{cvlist}
  \item[2023--present] Two PhD students (one as main advisor), five MSc theses.
  \item[2020--2023] Co-supervised three MSc theses; one led to a publication.
\end{cvlist}

\section{Invited talks}
\begin{cvlist}
  \item[Mar 2026] \emph{Robust Estimation of Sparse Graphs}. Statistics Seminar, Lakeshore University
  \item[Nov 2025] \emph{Learning on Networks under Contamination}. Workshop on High-Dimensional Statistics, Berlin
  \item[Jun 2025] \emph{Statistical Limits of Federated Graph Learning}. Keynote, Federated Learning Day, Springfield
  \item[Feb 2024] \emph{Adaptive Importance Sampling}. Probability Seminar, Riverside University
\end{cvlist}

\section{Professional service}
\subsection{Editorial and conference service}
\begin{cvlist}
  \item[2025--present] Associate Editor, \emph{Electronic Journal of Statistics}
  \item[2026] Program committee, ICML; area chair, AISTATS
  \item[2024] Co-organiser, Workshop on Statistics of Networks, Springfield (90 participants)
  \item[2020--present] Reviewer for JMLR, Annals of Statistics, Bernoulli, NeurIPS and others
\end{cvlist}
\subsection{Departmental service}
\begin{cvlist}
  \item[2024--present] Member, graduate admissions committee
  \item[2023--2024] Seminar organiser, Statistics Seminar
\end{cvlist}

% ---- Publications ----------------------------------------------------------------
% \nocite{*} prints every entry of refs.bib without citing it in the text.
% Each block prints one entry type. The label prefix (J, C, P) and the restarted
% numbering make an entry easy to refer to in a job application ("see J2").
\newcommand{\publist}[3]{%
  \begin{refcontext}[labelprefix=#1]
  \printbibliography[heading=pubs,title={#2},type=#3,resetnumbers=true]
  \end{refcontext}}

\section{Publications}
\nocite{*}
\publist{J}{Journal articles}{article}
\publist{C}{Conference papers (peer-reviewed)}{inproceedings}
\publist{P}{Preprints}{online}

\section{Languages and skills}
\begin{cvlist}
  \item[Languages] English (native), German (fluent), French (basic)
  \item[Software] R, Python, Julia; \LaTeX{} and Git for all papers and code
\end{cvlist}

\vspace{1em}
\textit{References are available on request.}

\end{document}
